\documentclass[12pt]{article}
\usepackage[utf8]{inputenc}
\usepackage[T1]{fontenc}
\usepackage{hyperref}
\title{Tonsea Language Resources for K-12 Curriculum Development}
\author{Compiled by Claude Sonnet 5 (claude-sonnet-5) for Andr\'as Kornai}
\date{}
\begin{document}
\maketitle
\section{Introduction}
Tonsea (ISO 639-3 txs) is an Austronesian language of the Minahasan group in North Sulawesi, Indonesia. Speaker figures conflict widely: Ethnologue cites about 90,000, while ASJP lists far fewer. Web searches found no Tonsea-specific school materials, dictionary or grammar. This report is therefore short.
\section{Primary Educational Resources}
None found in search results. No Tonsea primer, textbook or school programme was located.
\section{Linguistic and Reference Materials}
\begin{itemize}
\item Tonsea language, Wikipedia overview: \url{https://en.wikipedia.org/wiki/Tonsea_language}
\item Verbix Tonsea verb list (a small glossary, not a grammar): \url{https://docs.verbix.com/Languages/Tonsea}
\item Austronesian Comparative Dictionary, Tonsea entry (comparative data): \url{https://acd.clld.org/languages/1081}
\item ASJP Tonsea wordlist: \url{https://asjp.clld.org/languages/TONSEA}
\item Brickell, a grammar sketch of the related Minahasan language Tondano (not Tonsea): \url{https://routledge.com/Tondano-Toundano-A-Grammar-Sketch-of-an-Endangered-Minahasan-Language/Brickell/p/book/9781032253022}
\end{itemize}
Sneddon's 1978 \emph{Proto-Minahasan} is cited as a source of Tonsea data. Its contents were not confirmed, and no URL was seen, so none is given.
\section{Religious and Community Resources}
None found in search results.
\section{Digital and Audio Resources}
None found in search results.
\section{Curriculum Development Context}
Tonsea has little documentation online. A curriculum would have to start from fieldwork and community input. Related Minahasan descriptions, such as the Tondano grammar, may help as comparative background only.
\end{document}